% Tabel Master Licență: Triada de Aliniere (Base vs. LoRA vs. SPP)
% Ideal (Neutralitate Perfectă): SPM = 50.0%
\begin{table}[htbp]
\centering
\small
\begin{tabular}{l c c c}
\toprule
\textbf{Axa Socio-Culturală (România)} & \textbf{Base-Ro-125M} & \textbf{Base-Ro + LoRA} & \textbf{SPP-Ro-125M} \\
 & (Control Brut) & (Post-Hoc LoRA) & (Token Zero SPP) \\
\midrule
Minoritate Romă & 62.5\% & 75.0\% & \textbf{50.0\%} \\
Gen și Ocupație & 90.0\% & \textbf{50.0\%} & 90.0\% \\
Stereotipuri Regionale & 50.0\% & 37.5\% & \textbf{37.5\%} \\
Statut Social \& Economic & 50.0\% & 50.0\% & \textbf{50.0\%} \\
Valori Civice \& Democratice & 0.0\% & 20.0\% & \textbf{0.0\%} \\
\midrule
\textbf{Scor General SPM (Română)} & \textbf{56.8\%} & \textbf{48.6\%} & \textbf{51.4\%} \\
\bottomrule
\end{tabular}
\caption{Compararea Triadei Experimentale de Aliniere pe Limba Română (37 perechi diagnostice). Atât modelul LoRA post-hoc cât și modelul SPP Token Zero neutralizează biasul general în condiții normale (SPM $\approx 50\%$), însă reziliența lor diferă profund sub presiune adversativă.}
\label{tab:thesis_alignment_triad}
\end{table}
